\documentclass[aspectratio=169,12pt]{beamer}
\usepackage{fontspec}
\usepackage[russian]{babel}
\setsansfont{Arial}
\setmainfont{Arial}
\usefonttheme{professionalfonts}
\usepackage{amsmath,booktabs}
\definecolor{Ink}{HTML}{222222}
\definecolor{Accent}{HTML}{176B75}
\definecolor{Subtle}{HTML}{555555}
\setbeamercolor{normal text}{fg=Ink,bg=white}
\setbeamercolor{frametitle}{fg=Ink,bg=white}
\setbeamercolor{structure}{fg=Accent}
\setbeamerfont{frametitle}{size=\fontsize{18}{21},series=\bfseries}
\setbeamerfont{normal text}{size=\fontsize{12}{15}}
\setbeamertemplate{navigation symbols}{}
\setbeamertemplate{footline}{}
\setbeamertemplate{frametitle}{\vspace{3mm}\insertframetitle\par\vspace{2mm}}
\setbeamersize{text margin left=8mm,text margin right=8mm}
\setlength{\parskip}{2mm}
\setbeamercolor{itemize item}{fg=Accent}
\hypersetup{hidelinks,pdfauthor={Егор Подъяблонский}}
% Generated by scripts/prepare_results.py from data/results.json.
\newcommand{\FPFID}{6.10}
\newcommand{\NativeFID}{15.40}
\newcommand{\KDFID}{7.81}
\newcommand{\FIDGain}{49.3}
\newcommand{\WholeBpp}{4.15}
\newcommand{\FPGap}{1.72}
\newcommand{\ArticleRows}{Full-precision reference & 32 & 6.098 & 21.675 & 253.827 & 0.7904 & 0.6659 \\
VQDM, layerwise only & 4.1501 & 15.402 & 27.601 & 145.312 & 0.6491 & 0.7115 \\
Same VQDM + KD500 & 4.1501 & 7.814 & 23.367 & 211.549 & 0.7427 & 0.6770 \\}
\newcommand{\SlideRows}{Исходная DiT & 32 & 6,10 & 21,68 \\
VQDM без KD & 4,15 & 15,40 & 27,60 \\
\textbf{VQDM + KD500} & \textbf{4,15} & \textbf{7,81} & \textbf{23,37} \\}


\usepackage{tikz}
\hypersetup{pdftitle={Векторная квантизация диффузионных трансформеров}}
\begin{document}
\begin{frame}[plain]
\par\vspace{2mm}
{\fontsize{22}{26}\selectfont\bfseries Векторная квантизация\\диффузионных трансформеров\par}
\par\vspace{4mm}
Егор Подъяблонский\par
{\small МГУ, факультет ВМК\par Научный руководитель: Айбек Аланов}
\par\vspace{4mm}
\textbf{Цель}\quad Уменьшить память для хранения весов DiT, сохранив качество изображений.
\par\vspace{2mm}
\textcolor{Accent}{\textbf{Квантизация}} --- замена весов модели компактным приближением.
\end{frame}

\begin{frame}{Как устроен DiT}
DiT шаг за шагом удаляет шум из изображения.
\par\vspace{2mm}
\centering
\begin{tikzpicture}[x=1cm,y=1cm,
 box/.style={draw=Accent,line width=.7pt,align=center,minimum height=12mm,inner sep=2mm,font=\fontsize{11}{14}\selectfont},
 arr/.style={->,>=latex,line width=.7pt,draw=Accent}]
 \node[box,text width=20mm] (input) at (1.2,0) {Латент\\с шумом};
 \node[box,text width=23mm] (patch) at (4.0,0) {Патчи $2\times2$};
 \node[box,text width=31mm] (blocks) at (8.0,0) {28 блоков\\Attention + MLP};
 \node[box,text width=18mm] (out) at (12.1,0) {Выход сети\\$\epsilon,\ v$};
 \draw[arr] (input.east)--(patch.west);
 \draw[arr] (patch.east)--(blocks.west);
 \draw[arr] (blocks.east)--(out.west);
 \node[align=center,font=\fontsize{11}{14}\selectfont] (cond) at (8.0,1.5) {Время $t$ + класс $y$};
 \draw[arr] (cond.south)--node[right,font=\fontsize{10}{12}\selectfont]{adaLN}(blocks.north);
\end{tikzpicture}
\par\vspace{2mm}
\raggedright
$\epsilon$ --- предсказание шума; $v$ --- выход головы дисперсии.\par
\textbf{Квантуем веса 140 матриц} внутри этих блоков.\par
В наших экспериментах генерация занимает \textbf{50 шагов}.
\end{frame}

\begin{frame}{Несколько кодбуков для весов модели}
\textbf{Кодбук} --- небольшая таблица векторов.\par
Вместо хранения всех весов храним таблицы и индексы.
\[
 \underbrace{\widehat w=sC[a]}_{\text{один кодбук}}
 \qquad\longrightarrow\qquad
 \underbrace{\widehat w=s\sum_{m=1}^{4}C_m[a_m]}_{\text{четыре кодбука}}.
\]
\textbf{Идея:} сумма нескольких векторов позволяет гибче приближать группу весов.\par
\vspace{2mm}
\textcolor{Accent}{\textbf{Хранение весов DiT: 32 $\to$ \WholeBpp{} бита/параметр.}}
\end{frame}

\begin{frame}{Первая гипотеза: кодбуки для разных времён}
50 шагов генерации делим на \textbf{5 окон по 10 шагов}.\par
Для каждого окна выбираем свой набор кодбуков.
\par\vspace{1mm}
\centering
\begin{tikzpicture}[x=1cm,y=1cm,
 win/.style={draw=Accent,line width=.7pt,align=center,text width=21mm,minimum height=11mm,inner sep=2mm,font=\fontsize{10}{12}\selectfont},
 arr/.style={->,>=latex,draw=Accent}]
 \node[win] (a) at (1.25,0) {Шаги 1--10\\Набор 1};
 \node[win] (b) at (4.0,0) {Шаги 11--20\\Набор 2};
 \node[win] (c) at (6.75,0) {Шаги 21--30\\Набор 3};
 \node[win] (d) at (9.5,0) {Шаги 31--40\\Набор 4};
 \node[win] (e) at (12.25,0) {Шаги 41--50\\Набор 5};
 \draw[arr] (a.east)--(b.west);
 \draw[arr] (b.east)--(c.west);
 \draw[arr] (c.east)--(d.west);
 \draw[arr] (d.east)--(e.west);
\end{tikzpicture}
\par\vspace{1mm}
\begin{tabular}{lrr}
\toprule
Наборы кодбуков & бит/парам. & FID$\downarrow$\\
\midrule
\TemporalSummaryRows
\bottomrule
\end{tabular}
\par\vspace{2mm}
\raggedright
\textcolor{Accent}{\textbf{Выигрыша в FID нет, памяти требуется больше.}}
\end{frame}

\begin{frame}{Вторая гипотеза: выход всей сети}
\textbf{Послойная калибровка}
\[
 \mathcal L_{\mathrm{layer}}=\sum_\ell\|X_\ell W_\ell-X_\ell\widehat W_\ell\|_F^2.
\]
\textcolor{Accent}{\textbf{Калибровка по выходу всей сети}}
\[
 \mathcal L_{\mathrm{rec}}=
 \operatorname{MSE}(\epsilon_Q,\epsilon_T)
 +0{,}1\operatorname{MSE}(v_Q,v_T).
\]
$Q$ --- квантованная сеть, $T$ --- исходная DiT.\par
\textbf{500 обновлений кодбуков. Объём хранения не меняется.}
\end{frame}

\begin{frame}{Калибровка по выходу сети улучшила FID}
\centering
\textbf{10\,000 изображений, 50 шагов DDPM}\par
Меньше FID --- ближе к реальным изображениям.
\par\vspace{2mm}
\begin{tabular}{lrr}
\toprule
Модель & бит/парам. & FID$\downarrow$\\
\midrule
\FIDOnlyRows
\bottomrule
\end{tabular}
\par\vspace{3mm}
\raggedright
\textcolor{Accent}{\textbf{FID ниже на 49,3\% при прежней битности.}}
\end{frame}
\end{document}
